%% notes-tex/031-yeast-chemical-space/sections/1-question.tex

\section{The question\secstatus{todo}}
\label{sec:question}

Experiment 031 represents a dosed inhibitor, a medium component and a metabolic-model
metabolite in one embedding space, which works only to the extent that the molecules are
actually named and structured. That raised a broader question worth scoping before it is
acted on: how many distinct chemical species does a yeast cell contain, how many does the
genome-scale reconstruction name, and could a retrobiosynthetic method enumerate the rest.

The motivation is a physical accounting rather than a modeling convenience. If every species
carried a resolved structure with an atom mapping, then reaction rules, thermodynamics and
molecular embeddings would all read the same object, and gap filling would become a question
about chemistry rather than about identifiers. The practical target is the set of
intermediates reachable from a defined medium, including the uncatalyzed and spontaneous
chemistry that constraint-based reconstruction tends to leave out.

This note measures what the reconstruction gives us, reports what the literature has already
settled, and states what is left. It is short on purpose: the point is to decide whether the
side quest is worth starting.
